% 图 4-1 KSE SoC 结构框图
\documentclass[border=4pt]{standalone}
% 所有框图共用的样式（standalone + TikZ，xelatex/tectonic 编译）
\usepackage{fontspec}
\usepackage{xeCJK}
\setCJKmainfont{Noto Sans CJK SC}[BoldFont=Noto Sans CJK SC Bold]
\setCJKsansfont{Noto Sans CJK SC}[BoldFont=Noto Sans CJK SC Bold]
\renewcommand{\familydefault}{\sfdefault}
\usepackage{tikz}
\usetikzlibrary{positioning,arrows.meta,fit,calc,backgrounds,shapes.geometric,decorations.pathreplacing,patterns}
% 统一配色：灰蓝主色，橙色强调，紫色辅助
\definecolor{cBlue}{HTML}{2E6DA4}
\definecolor{cBlueL}{HTML}{DCE7F2}
\definecolor{cBlueM}{HTML}{A9C3DD}
\definecolor{cOrange}{HTML}{C9732E}
\definecolor{cOrangeL}{HTML}{F6E3D2}
\definecolor{cPurple}{HTML}{8A5BB0}
\definecolor{cPurpleL}{HTML}{ECE3F4}
\definecolor{cGray}{HTML}{5B6670}
\definecolor{cGrayL}{HTML}{EEF0F2}
\definecolor{cInk}{HTML}{25303B}
\tikzset{
  >={Stealth[length=2.2mm,width=1.6mm]},
  every node/.style={font=\small, text=cInk},
  blk/.style={draw=cBlue, fill=cBlueL, line width=0.6pt, rounded corners=1.5pt, align=center, minimum height=7mm, inner sep=3pt},
  blkO/.style={blk, draw=cOrange, fill=cOrangeL},
  blkP/.style={blk, draw=cPurple, fill=cPurpleL},
  blkG/.style={blk, draw=cGray, fill=cGrayL},
  blkW/.style={blk, draw=cGray, fill=white},
  grp/.style={draw=cGray, dashed, line width=0.5pt, rounded corners=3pt, inner sep=6pt},
  grpS/.style={draw=cBlue, line width=0.8pt, rounded corners=3pt, inner sep=6pt},
  lbl/.style={font=\footnotesize, text=cGray, inner sep=1pt},
  ttl/.style={font=\small\bfseries, text=cInk, inner sep=1pt},
  ar/.style={->, draw=cInk, line width=0.6pt},
  arB/.style={->, draw=cBlue, line width=0.9pt},
  arO/.style={->, draw=cOrange, line width=0.9pt},
  arD/.style={->, draw=cGray, line width=0.6pt, dashed},
  bus/.style={draw=cGray, line width=2.4pt},
}

\begin{document}
\begin{tikzpicture}[
  u/.style={blk, font=\footnotesize, minimum height=7mm, inner sep=2pt},
  uo/.style={blkO, font=\footnotesize, minimum height=7mm, inner sep=2pt},
  ug/.style={blkG, font=\footnotesize, minimum height=7mm, inner sep=2pt},
  up/.style={blkP, font=\footnotesize, minimum height=7mm, inner sep=2pt},
]
\node[ug, minimum width=4.2cm, minimum height=1.2cm] (cpu) at (3.5,0) {\textbf{E203}（RV32IMAC）\\ITCM 64 KiB / DTCM 64 KiB};
\node[ug, minimum width=4.2cm] (biu) at (3.5,-1.3) {BIU};
\draw[ar] (cpu) -- (biu);
\node[ug, minimum width=2.0cm, minimum height=1.0cm] (plic) at (11.0,-1.3) {PLIC\\\scriptsize 0x0C00\_0000};
\draw[arD] (plic.north) |- node[lbl, pos=0.75, above] {外部中断} (cpu.east);
% 加速器子系统
\node[up, minimum width=6.2cm] (xbar) at (2.6,-3.6) {系统总线仲裁器（3 进 1 出，DMA > 加载器 > CPU）};
\node[u, minimum width=2.9cm, minimum height=1.25cm] (bk0) at (1.05,-5.1) {SRAM 控制器\\\scriptsize 64 × RAMSP4096X16};
\node[up, minimum width=2.9cm, minimum height=1.25cm] (bk1) at (4.15,-5.1) {进度寄存器\\\scriptsize BASE / WPTR / DONE};
\node[lbl, anchor=north] (sl) at (2.6,-5.85) {共享 SRAM 512 KiB（0x4000\_0000），单个 32 位从口；门铃 0x4008\_0000};
\draw[<->, draw=cInk, line width=0.6pt] (bk0.north) -- (bk0.north |- xbar.south);
\draw[<->, draw=cInk, line width=0.6pt] (bk1.north) -- (bk1.north |- xbar.south);
\node[uo, minimum width=3.1cm, minimum height=1.6cm, align=center] (wrap) at (8.6,-3.6) {\textbf{引擎包装}\\\scriptsize MMIO 寄存器\\\scriptsize 命令队列 / 参数段队列\\\scriptsize 流式 DMA};
\node[u, minimum width=3.1cm, minimum height=1.25cm, draw=cBlue, fill=cBlueM!50] (eng) at (8.6,-5.4) {\textbf{引擎本体}（kse\_top）};
\draw[<->, draw=cInk, line width=0.6pt] (wrap) -- (eng);
\draw[ar] (wrap.west) -- node[lbl, above, pos=0.5] {DMA} (xbar.east);
\coordinate (xtop) at (2.6,-2.75);
\begin{scope}[on background layer]
\node[draw=cGray, dashed, line width=0.6pt, rounded corners=3pt, fit=(xbar)(bk0)(bk1)(sl)(wrap)(eng)(xtop), inner sep=7pt] (sys) {};
\end{scope}
\node[ttl, anchor=north west] at ($(sys.south west)+(0,-0.05)$) {\footnotesize 加速器子系统（hfclk 150 MHz，单时钟域）};
% 总线
\draw[ar] (2.1,-1.65) -- node[lbl, left, pos=0.25] {存储总线 O5（CPU）} (2.1,-3.25);
\draw[ar] (4.9,-1.65) -- (4.9,-2.2) -- node[lbl, above] {外设总线 PPI O12（MMIO 0x1100\_0000）} (8.6,-2.2) -- (wrap.north);
\draw[ar] (biu.east) -- (plic.west);
\draw[ar, draw=cOrange, line width=0.9pt] (wrap.east |- 0,-3.3) -- (11.0,-3.3) -- node[lbl, right, pos=0.5, text=cOrange!80!black] {中断源 17} (plic.south);
% 加载器
\node[ug, minimum width=2.4cm, minimum height=1.0cm] (ldr) at (-3.4,-3.6) {参数加载器\\\scriptsize 32 位 ICB 写入};
\draw[ar] (ldr.east) -- node[lbl, below, pos=0.45, align=center] {参数\\门铃} (xbar.west);
\draw[arD] (bk1.east) -- (6.3,-5.1) -- node[lbl, pos=0.35, right] {放行上界} (6.3,-3.95) -- (wrap.west |- 0,-3.95);
\end{tikzpicture}
\end{document}
